\scanpage{31}
\mspage{11}

d'où des coefficients dans les $\mathbb Q_\ell$, et un premier problème est s'ils sont dans $\mathbb Q$. On sait (si $X$ est projectif non singulier) ramener ce pb à un pb de cycles algébriques du type Lefschetz.

Mais --- c'est là le ``\unsure{point}'' que je désire faire --- $\overline{\mathrm m}$ si ce pb ``\unsure{cru}'' \inserted{i.e. ``\unsure{non tordu}''} est résolu par l'affirmative, on \inserted{ne} pourra encore conclure que les coeff. sont des entiers, mais seulement \inserted{de $\mathbb Z[1/p]$, i.e.} de \deleted{entiers $\ell$-adiques} la forme $n/p^{\unsure{\delta}}$, avec $n\in\mathbb Z$. Ceci provient de la restriction
\[
\ell\ne\operatorname{car} k,
\]
de ce \uline{qu'il manque une théorie cohomologique $p$-adique pour les variétés algébriques en car. $p$}. Faute d'avoir une théorie de ce type, on ne saura

\margin{30}
